%=====================================================================
%  I. PENDAHULUAN
%=====================================================================
\section{PENDAHULUAN}

\lettrine[lines=2, findent=2pt, nindent=0pt]{I}{saac} Newton memandang gravitasi
sebagai gaya mekanis yang menarik dua benda bermassa satu sama lain. Einstein
mengubah pandangan tersebut secara mendasar: gravitasi bukanlah gaya, melainkan
manifestasi geometris dari kelengkungan ruang-waktu yang disebabkan oleh
distribusi materi dan energi \cite{einstein1916}. Terdapat hubungan timbal balik
di dalamnya, yakni materi dan energi menentukan bagaimana ruang-waktu melengkung,
dan sebaliknya kelengkungan ruang-waktu menentukan bagaimana materi bergerak
melaluinya. Hubungan tersebut dirumuskan secara matematis melalui persamaan medan
Einstein, yang mengaitkan tensor Einstein sebagai representasi geometri dengan
tensor energi-momentum sebagai representasi kandungan materi \cite{wald1984}.

Persamaan medan Einstein bersifat sangat umum dan tidak mudah diselesaikan secara
langsung untuk sembarang distribusi materi. Kasus paling sederhana yang dapat
ditinjau adalah kondisi vakum, yaitu ketika tidak terdapat materi maupun energi di
luar suatu objek masif, sehingga ruas kanan persamaan medan Einstein lenyap.
Penyelesaian persamaan medan Einstein pada kondisi vakum dengan asumsi ruang-waktu
statis dan bersimetri bola inilah yang pertama kali diperoleh sebagai solusi eksak,
yaitu solusi Schwarzschild \cite{schwarzschild1916}. Solusi ini menggambarkan
geometri ruang-waktu di sekitar objek masif bersimetri bola melalui sebuah fungsi
metrik tunggal yang hanya bergantung pada massa objek \cite{carroll2004}.

Solusi Schwarzschild tidak hanya mendeskripsikan geometri di luar suatu bintang,
tetapi juga mendeskripsikan lubang hitam paling sederhana, yakni lubang hitam yang
tidak berotasi dan tidak bermuatan. Lubang hitam merupakan daerah ruang-waktu
dengan medan gravitasi sedemikian kuat sehingga tidak ada benda apa pun, termasuk
cahaya, yang dapat lepas dari dalamnya setelah melewati suatu batas tertentu
\cite{frolov1998, misner1973}. Batas tersebut dikenal sebagai horizon peristiwa,
yaitu permukaan yang memisahkan daerah yang masih dapat berkomunikasi dengan
pengamat di kejauhan dari daerah di baliknya yang tidak lagi dapat melakukannya.
Berbeda dengan horizon peristiwa yang semata-mata merupakan batas kausal,
singularitas fisis yang sesungguhnya terletak di pusat lubang hitam, tempat
kelengkungan ruang-waktu menjadi tak hingga dan deskripsi geometri klasik tidak
lagi berlaku \cite{frolov1998}.

Karena sifat ekstremnya, lubang hitam menjadi sistem yang ideal untuk mempelajari
karakteristik ruang-waktu secara lebih mendalam. Salah satu cara paling alami untuk
menyelidiki karakteristik tersebut adalah dengan memberikan gangguan kecil pada
sistem, lalu meninjau bagaimana ruang-waktu meresponsnya. Pendekatan inilah yang
dikenal sebagai teori perturbasi lubang hitam. Regge dan Wheeler (1957) mengawali
pendekatan tersebut dengan menuliskan metrik yang terganggu sebagai jumlah dari
metrik latar belakang dan suku gangguan kecil, kemudian menunjukkan bahwa gangguan
kecil pada lubang hitam Schwarzschild menghasilkan sebuah persamaan gelombang
dengan potensial efektif \cite{regge1957}. Temuan ini membuka perspektif bahwa
respons dinamis lubang hitam terhadap gangguan dapat dipelajari secara sistematis
melalui analisis perturbasi linier \cite{vishveshwara1970}.

Gangguan pada lubang hitam dapat dibedakan berdasarkan spin medan yang meninjaunya,
yaitu medan skalar (spin $s = 0$), medan elektromagnetik ($s = 1$), dan medan
gravitasi ($s = 2$) \cite{konoplya2011}. Ketiganya memiliki tingkat kerumitan
analisis yang berbeda. Medan skalar merupakan jenis gangguan yang paling sederhana
karena dinamikanya tidak memberikan pengaruh balik terhadap geometri ruang-waktu
latar belakang, sehingga medan skalar dapat diperlakukan sebagai medan uji.
Kesederhanaan inilah yang menjadikan medan skalar sebagai titik awal yang wajar
sebelum meninjau jenis gangguan yang lebih kompleks, sekaligus tetap
mempertahankan seluruh struktur matematis pokok dari persoalan perturbasi lubang
hitam \cite{chandrasekhar1983, berti2016}.

Ciri utama dari struktur matematis tersebut adalah bahwa persamaan gerak medan uji
pada latar belakang statis dan bersimetri bola dapat direduksi menjadi sebuah
persamaan diferensial biasa berbentuk persamaan gelombang satu dimensi. Reduksi ini
dimungkinkan karena simetri latar belakang memungkinkan pemisahan variabel sudut
melalui basis harmonik sferis, sedangkan sifat statis memungkinkan pemisahan
variabel waktu melalui bergantungan harmonik \cite{schutz2009}. Bagian radial yang
tersisa kemudian dapat diubah bentuknya menjadi persamaan yang menyerupai persamaan
Schr\"{o}dinger tak gayut waktu melalui pengenalan koordinat \textit{tortoise},
sehingga seluruh pengaruh geometri lubang hitam terangkum di dalam sebuah potensial
efektif \cite{konoplya2011, kokkotas1999}.

Potensial efektif inilah yang menjadi objek fisis pokok dalam teori perturbasi
lubang hitam. Potensial tersebut berperan sebagai penghalang (\textit{barrier})
yang menghamburkan gelombang yang datang dari horizon peristiwa maupun dari jarak
jauh, sehingga bentuk, ketinggian, dan letak puncaknya menentukan bagaimana energi
gangguan ditransmisikan, dipantulkan, dan pada akhirnya diluruhkan oleh lubang
hitam \cite{nollert1999}. Dengan demikian, karakterisasi potensial efektif merupakan
langkah yang harus dituntaskan lebih dahulu sebelum kajian dinamis apa pun terhadap
respons lubang hitam dapat dilakukan \cite{teukolsky1973}.

Berdasarkan uraian tersebut, penelitian ini bertujuan untuk menurunkan persamaan
perturbasi medan skalar pada ruang-waktu Schwarzschild secara analitik, mulai dari
persamaan Klein-Gordon dalam ruang-waktu lengkung hingga diperoleh bentuk persamaan
master beserta potensial efektifnya, serta untuk menganalisis karakteristik
potensial efektif tersebut terhadap variasi bilangan momentum sudut, massa lubang
hitam, dan massa medan skalar itu sendiri.
